\documentclass[10pt,a4paper]{amsart}
\usepackage[margin=19mm]{geometry}
\usepackage{amsmath,amssymb,mathtools}
\usepackage[scr=rsfs,cal=euler]{mathalfa}
\usepackage{xfrac}
\usepackage[unicode,hidelinks,bookmarks=false]{hyperref}
\newcommand{\ie}{i.e.\ }
\newcommand{\cf}{cf.\ }
\newcommand{\resp}{respectively\ }
\newcommand{\Subsection}{\subsection}
\providecommand{\nobreakdash}{\nobreak\hbox{-}}
\setlength{\emergencystretch}{2em}
\multlinegap0pt
\usepackage{amssymb}
\newtheorem{lemma}{Lemma}
\DeclareMathOperator{\End}{End}
\title{Rota--Baxter Operators on Truncated Polynomial Algebras}
\author{Azhar Farooq}
\date{Source: arXiv:2605.15670v1 (CC BY 4.0)}
\begin{document}
\maketitle

\noindent\emph{Source and licence.} Rota--Baxter Operators on Truncated Polynomial Algebras. Version arXiv:2605.15670v1 (CC BY 4.0), distributed by arXiv under the Creative Commons Attribution 4.0 International licence, http://creativecommons.org/licenses/by/4.0/, as recorded in the arXiv licence metadata for this exact version and shown on the abstract page https://arxiv.org/abs/2605.15670v1. Full licence notice: \texttt{licenses/arxiv-2605.15670-CC-BY-4.0.txt}.\par\smallskip

\noindent\textit{Editorial scope.} Counted unit: the article's lemma lem:wt1\_form (source0.tex lines 231--233). The article's complete original proof of it is delivered with it (source0.tex lines 235--284, one \texttt{proof} environment of 50 lines). No mathematics is written, repaired or re-derived: the statement and the proof are the article's own lines, in the article's own notation, and no retained line is substituted or re-flowed.  Retained uncounted as the source-local context the proof needs: lines 59--61.  Nothing was omitted from any retained span, so the package carries no omission record.  No retained line contains a citation, so the wrapper carries no bibliography and no bibliography entry.  No retained line contains a figure environment, an image-inclusion command or drawing code, so no image is delivered and the package declares no asset dependency.  Editorial formatting only, in addition: an AMS \texttt{amsart} preamble that restates the article's own declarations and macros used by the retained lines, namely the environments \texttt{lemma} on separate counters, together with the standard packages those lines need.  The retained environments print in this document as Lemma 1 in order of appearance, whereas the article prints the same items with the numbering of its own full document.  The complete native source of the article as distributed in the arXiv e-print for this exact version is bundled unchanged as \texttt{sources/200670/source0.tex}.
\begin{equation}\label{RBdef}
    P(x)P(y) \;=\; P\bigl(xP(y)+P(x)y + \lambda\,xy\bigr)
\end{equation}

\begin{lemma}\label{lem:wt1_form}
If $P : R \to R$ is a Rota--Baxter operator of weight one, then $P|_V$ maps $V$ to $V$, and  $P(1) \in \{0, -1\}$.  Moreover, for all $u \in V$, $P(u) = L(u)$ with $L \in \End_K(V)$ satisfying $L^2 + L = 0$.
\end{lemma}

\begin{proof}
Write $P(1) = \alpha + v_0$ with $\alpha \in K$, $v_0 \in V$, and for $u \in V$ write $P(u) = f(u) + L(u)$ where $f : V \to K$ is linear and $L \in \End_K(V)$.

First take $x, y \in V$.  The weight‑one identity \eqref{RBdef} with $\lambda=1$ gives
\[
P(x)P(y) = P\bigl(x P(y) + P(x) y + xy\bigr).
\]
Since $V^2=0$, we have $xy=0$, $x L(y)=0$, and $L(x) y=0$.  Then
\begin{align*}
\text{LHS} &= (f(x)+L(x))(f(y)+L(y)) = f(x)f(y) + f(x)L(y) + f(y)L(x),\\
\text{RHS} &= P\bigl(f(y)x + f(x)y\bigr) = f(y)P(x) + f(x)P(y)\\
 &= 2f(x)f(y) + f(y)L(x) + f(x)L(y).
\end{align*}
Equating scalar parts yields $f(x)f(y) = 2f(x)f(y)$, hence $f(x)f(y)=0$ for all $x,y$.  A non‑zero linear functional would have $f(u_0)=1$ for some $u_0$, then $f(u_0)^2=1$, contradiction.  Thus $f=0$ and $P(V) \subseteq V$.

Now set $x = y = 1$ in \eqref{RBdef}:
\[
P(1)^2 = P(2P(1) + 1).
\]
Using $P(1)=\alpha+v_0$ and $P|_V = L$,
\begin{align*}
\text{LHS} &= \alpha^2 + 2\alpha v_0,\\
\text{RHS} &= P(2\alpha + 2v_0 + 1) = (2\alpha+1)P(1) + 2P(v_0) \\
&= (2\alpha+1)(\alpha+v_0) + 2L(v_0) = 2\alpha^2+\alpha + (2\alpha+1)v_0 + 2L(v_0).
\end{align*}
Comparing components gives
\begin{align*}
\alpha^2 &= 2\alpha^2 + \alpha \;\Longrightarrow\; \alpha(\alpha+1)=0,\\
2\alpha v_0 &= (2\alpha+1)v_0 + 2L(v_0) \;\Longrightarrow\; L(v_0) = -\tfrac12 v_0. \tag{B}
\end{align*}
So $\alpha \in \{0, -1\}$.

Next take $x = u \in V$, $y = 1$:
\[
P(u)P(1) = P\bigl(u P(1) + P(u) + u\bigr).
\]
Compute:
\begin{align*}
\text{LHS} &= L(u)(\alpha+v_0) = \alpha L(u),\\
\text{RHS} &= P\bigl(u(\alpha+v_0) + L(u) + u\bigr) = P\bigl(\alpha u + L(u) + u\bigr)\\
&= (\alpha+1)L(u) + L^2(u).
\end{align*}
Hence $\alpha L(u) = (\alpha+1)L(u) + L^2(u)$, i.e. $L^2 + L = 0$ on $V$.

Finally, apply $L$ to Equation (B):
\[
L^2(v_0) = -\tfrac12 L(v_0) = \tfrac14 v_0.
\]
But $L^2 = -L$, so $L^2(v_0) = -L(v_0) = \tfrac12 v_0$.  Thus $\tfrac14 v_0 = \tfrac12 v_0$, forcing $v_0=0$.  Hence $P(1) = \alpha \in \{0,-1\}$.
\end{proof}

\end{document}
